\documentclass[10pt]{article}
\usepackage[T1]{fontenc}
\usepackage{amsmath,amssymb,amsthm,geometry,hyperref}
\geometry{margin=1in}
\newcommand{\T}{\mathbb T}
\newcommand{\E}{\mathbb E}
\newcommand{\Pp}{\mathbb P}
\title{Finite-particle singular realization and heat passage}
\author{HOCF--RIESZ, Round 006}
\date{18 September 2026}
\begin{document}
\maketitle
\noindent\textbf{Scope and status.}
This memorandum records the finite-$N$ assertion THM--026. Its complete
construction, fresh reconstruction and hostile review have separate immutable
Markdown reports; subsequent dispositions are in the canonical theorem ledger.
The result concerns particle realization and law domination. It does not supply
the domain of a singular pair corrector or estimates uniform in particle number.

\section*{Model and precise conclusions}
The torus has Haar mass one and characters $e^{2\pi i k\cdot x}$. Fix
$d\ge3$, $0<s\le d-2$, $N\ge2$, $0\le\nu<\infty$, a finite horizon $T$,
and a smooth real periodic potential $V$. Set $b=-\nabla V$, $K=-\nabla g_s$,
where
\[
 \widehat g_s(0)=0,\qquad
 \widehat g_s(k)=\pi^{s-d/2}
 \frac{\Gamma((d-s)/2)}{\Gamma(s/2)}|k|^{s-d},\quad k\ne0.
\]
The previously proved heat representation gives $g_s(z)=|z|^{-s}+h_s(z)$
with a smooth even local remainder, $g_s,K\in L^1$, and
\[
 D:=\operatorname{div}K=
 \begin{cases}
 s(d-2-s)g_{s+2}\,dz,&s<d-2,\\
 c_d(\delta_0-dz),&s=d-2,
 \end{cases}
 \qquad c_d=(d-2)|\mathbb S^{d-1}|.
\]
Choose $a,\kappa\ge0$ with $\operatorname{div}b\ge-a$ and
$D\ge-\kappa\,dz$, and put $C_N=Na+(N-1)\kappa$.
With independent standard Brownian motions, the equation is
\[
 dX_i=\left[b(X_i)+\frac1N\sum_{j\ne i}K(X_i-X_j)\right]dt
       +\sqrt{2\nu}\,dW_i.
\]
From every fixed distinct configuration it has a unique global collision-free
pathwise realization. The path and transition kernels are jointly Borel and
have the deterministic-time conditional Markov property. Each fixed realized
finite horizon has a positive minimum pair distance. No single exceptional
set for all uncountably many starts is asserted.

For the same initial state and Brownian motions, replacing $K$ by its positive
heat convolution $K_\varepsilon$ gives
\[
 \sup_{0\le t\le T}\max_i
 \operatorname{dist}_{\T^d}(X_i^\varepsilon(t),X_i(t))\longrightarrow0
 \quad\text{almost surely as }\varepsilon\downarrow0.
\]
This is a fixed-parameter full-family limit, with no uniform rate. An initial
probability density $F_0\in L^\infty((\T^d)^N)$, independent of the Brownian
motions, evolves to a density with
\[
 F_t\le\|F_0\|_\infty e^{C_Nt}\quad\text{almost everywhere}.
\]
For iid initial density bounded by $M$, the factor is $M^Ne^{C_Nt}$.
The estimate does not propagate independence or give a uniform bound in $N$.

\section*{Energy controls every partial collision}
Write $v_*=\min V$, $g_*=\inf_{z\ne0}g_s(z)>-\infty$, and define
\[
 H_N=\sum_iV(x_i)+\frac1N\sum_{i<j}g_s(x_i-x_j),\qquad
 \mathcal E_N=H_N-Nv_*-\frac{N-1}{2}g_*.
\]
Each summand in
\[
 \mathcal E_N=\sum_i(V(x_i)-v_*)+
        \frac1N\sum_{i<j}(g_s(x_i-x_j)-g_*)
\]
is nonnegative. Thus a single colliding pair forces the energy to infinity,
including at simultaneous clusters or disjoint partial collisions. All energy
sublevels are compact in the collision-free configuration space. If
$\delta=\min_{i<j}\operatorname{dist}(x_i,x_j)\le r_0=1/4$, a fixed
$A>0$ determined by the local remainder gives
\[
 \mathcal E_N\ge\frac{\delta^{-s}-A}{N}.
\]
The complete drift $B$ is $-\nabla H_N$ and the full trace is
\[
 \Delta_{Nd}H_N=\sum_i\Delta V(x_i)+\frac2N\sum_{i<j}\Delta g_s(x_i-x_j)
 \le C_N.
\]
The coefficient is $2/N$. At Coulomb, the punctured periodic Laplacian is
$\Delta g_s=c_d$; it is not the zero Euclidean principal Laplacian.
For $Q_N=C_N-\Delta H_N\ge0$, the exact generator identity is
\[
 L_N\mathcal E_N=\nu C_N-|B|^2-\nu Q_N.
\]
The square is the full squared total drift. Its triple cross terms can have
either sign and are not discarded. Three equally spaced collinear particles
have zero leading force on the middle particle but divergent total pair energy;
no individual-force coercivity is needed.

\section*{Construction, stopping and global integrability}
Smooth periodic force cutoffs agree with $K$ away from a shrinking collision
tube. Subtracting the continuous additive driving path gives a globally
Lipschitz integral equation; Picard iteration and Gronwall construct its unique
solution and Borel dependence. Local uniqueness patches these solutions up to
collision. On the torus the only possible finite lifetime approaches that set.

Let $\sigma_R$ be the first energy-level exit. Before this stop the coefficients
and energy derivatives lie in a fixed compact smooth region. Smooth It\^o
calculus therefore gives
\begin{align*}
 \mathcal E_N(X_{t\wedge\sigma_R})
 &+\int_0^{t\wedge\sigma_R}|B|^2\,dr
 +\nu\int_0^{t\wedge\sigma_R}Q_N\,dr
 =\mathcal E_N(x)+\nu C_N(t\wedge\sigma_R)+M_R(t),\\
 [M_R]_t&=2\nu\int_0^{t\wedge\sigma_R}|B|^2\,dr.
\end{align*}
The stopped integrand is bounded, so $M_R$ is a true square-integrable
martingale. Taking expectations, using nonnegativity, yields
\[
 \Pp_x(\sigma_R\le T)\le\frac{\mathcal E_N(x)+\nu C_NT}{R}.
\]
Energy coercivity implies that these stops exhaust the local lifetime.
Letting $R$ tend to infinity proves noncollision and global uniqueness, before
using positive minimum path separation.

For independent initial data with $J_0=\E\mathcal E_N(X_0)<\infty$, Fatou
first proves integrability of $\int_0^T|B|^2$ and $\int_0^T\nu Q_N$.
The unstopped energy martingale is then square integrable. Its omitted bracket
has expectation tending to zero, so the stopped martingales converge in $L^2$.
Passing the localized identity and taking expectations now gives the equality
\[
 \E\mathcal E_N(X_t)+\E\int_0^t|B|^2\,dr
        +\E\int_0^t\nu Q_N\,dr=J_0+\nu C_Nt.
\]
At $\nu=0$ the last integral is identically zero; no division by $\nu$ occurs.
The same stopped argument at a distance threshold $\rho$ gives
\[
 \Pp\left(\min_{[0,T]}\delta(X_t)\le\rho\right)
 \le\frac{N(J_0+\nu C_NT)}{\rho^{-s}-A}
 \le2N(J_0+\nu C_NT)\rho^s
\]
when $0<\rho\le r_0$ and $\rho^{-s}\ge2A$.
For bounded $F_0$, the required energy mean is finite because
\[
 \int\mathcal E_N\,dX=N\left(\int V-v_*\right)-\frac{N-1}{2}g_*<\infty.
\]
This initial integrability is established before any density propagation.

The countable cutoff construction yields a Borel set of good start/path pairs,
whose section has probability one at each distinct start. Independent Brownian
increments and Fubini at a random current state justify deterministic-time
restart and the Markov property. A universal exceptional set is unnecessary.

\section*{Heat coupling and density passage}
On every compact set away from zero, $K_\varepsilon\to K$ in $C^1$.
To prove this, separate a smooth part agreeing with $K$ near that set from an
$L^1$ distant part. Gaussian derivative tails bound the latter by an inverse
power of $\varepsilon$ times $e^{-c/\varepsilon}$. Fix a noncolliding path and
its positive minimum separation. The Brownian terms cancel in the same-noise
difference equation. A stopped Gronwall bound prevents the smooth trajectory
from leaving a small neighborhood of that path for every sufficiently small
$\varepsilon$. This proves the full-family uniform path limit.

The regularized full drift has exact divergence
\[
 \operatorname{div}_{Nd}B^\varepsilon
 =\sum_i\operatorname{div}b(x_i)+\frac2N\sum_{i<j}D_\varepsilon(x_i-x_j)
 \ge-C_N.
\]
For each continuous driving path its smooth flow is a spatial diffeomorphism:
solve the translated integral equation backwards for its inverse. Its ordinary
variational equation gives a Jacobian determinant at least $e^{-C_Nt}$.
Change of variables, the bound on $F_0$, and Brownian expectation give
\[
 \E G(X_t^\varepsilon)\le\|F_0\|_\infty e^{C_Nt}\int G\,dX
 \quad(G\ge0\text{ continuous}).
\]
The same-noise convergence and bounded convergence first pass this inequality
for continuous tests. Increasing continuous approximations of open-set
indicators, followed by Haar outer regularity, extend it to all Borel sets.
This proves the asserted density domination without assuming a singular
Fokker--Planck solution, convergence for all Borel terminal tests, or an inverse
for the limiting singular flow.

\section*{The remaining corrector-domain line}
These results do not justify applying It\^o calculus to the bounded Borel
full pair inverse from Round005. Its spatial and time derivatives,
differentiated background contractions, singular generator identity, and
residual/bracket approximation remain separate obligations. A finite-$N$
density bound transfers an already proved integrable function; it cannot
supply a missing derivative or turn a nonintegrable expression into an
integrable one. Uniform evolved-law estimates and the critical hierarchy
remain open, with the original energy-floor and microscopic regimes unchanged.
\end{document}
